% Appendix - ICML style

\section{Dataset Details}
\label{app:data}

\paragraph{Porto Taxi.}
We use the Porto Taxi dataset \citep{porto_taxi_dataset}, comprising $\sim$108K training routes collected from taxi GPS traces in Porto, Portugal.
The road network is extracted from OpenStreetMap using OSMnx, yielding $\sim$35K unique ways.
Node coordinates are projected to a $1024 \times 1024$ grid system.

\paragraph{Way-level tokenization.}
Ground-truth GPS trajectories are map-matched to the road graph.
Each matched node sequence is converted to a way sequence by identifying the OSM way containing each node, then performing consecutive deduplication (merging repeated adjacent ways into a single token).
Routes with excessive map-matching failures are filtered.

\paragraph{OD-disjoint split.}
To evaluate generalization to unseen origin-destination pairs, we partition routes such that no OD pair appears in both training and test sets (dev10p split).
All reported metrics use 5{,}000 test routes.

\section{Implementation Details}
\label{app:impl}

\paragraph{Autoencoder.}
The way encoder embeds each way as a 256-dimensional token using learned embeddings for way identity (vocabulary $\sim$35K), road tier (3 classes: major/minor/service), and positional features (grid coordinates, segment length).
The Perceiver compressor uses 8 learnable query tokens with 4 cross-attention layers and 4 attention heads, producing $8 \times 256 = 2048$-dimensional latent representations.

The $\beta$-VAE bottleneck projects 2048 dimensions to $D = 64$ via MLP, applies reparameterization, and projects back to 2048 via a decoder MLP.
$\beta$ is warmed up linearly from 0 to 0.01 over the first 30 epochs.
The decoder is a graph-constrained autoregressive model with cross-attention to the 8 latent tokens, candidate-aware query mechanism (max 32 candidates per step), and past-context encoding (last 16 steps).

Training uses AdamW ($\mathrm{lr} = 5 \times 10^{-4}$, weight decay $10^{-4}$) with cross-entropy loss on the way prediction task, for 100 epochs with batch size 128.

\paragraph{Flow model.}
The conditional flow model is a 6-layer transformer with $d_{\mathrm{model}} = 256$ and 4 attention heads.
Conditioning features (origin embedding, destination embedding, time features via sinusoidal encoding) are injected through cross-attention.
The flow target is $\boldsymbol{\mu} \in \mathbb{R}^{64}$, the posterior mean from the trained $\beta$-VAE.
Training uses the conditional flow matching objective \citep{lipman2023flow} with AdamW ($\mathrm{lr} = 10^{-4}$) for 200 epochs.

\paragraph{Inference.}
At inference, the flow model generates $K = 16$ samples $\{\boldsymbol{\mu}_k\}$ using 100 ODE integration steps.
Each $\boldsymbol{\mu}_k$ is mapped to $8 \times 256$ latent tokens via the bottleneck decoder MLP, then decoded autoregressively with anti-loop penalty ($k = 4$, $\alpha = 2.0$).
The maximum decoding length is 200 way tokens.

\paragraph{AR baselines.}
The RNN baseline is a 2-layer GRU with hidden size 256, trained to predict the next way given the current way and conditioning.
The Transformer baseline is a 4-layer causal Transformer with $d_{\mathrm{model}} = 256$.
Both use beam search with $K = 10$ at inference.
Both are trained on the same way-level tokenized routes with cross-entropy loss.

\section{Coverage-$\tau$ AUC Computation}
\label{app:covauc}

For each OD group with $N_{\mathrm{gt}} \geq 2$ ground-truth routes, we compute:
\begin{equation}
  \mathrm{Coverage}@K(\tau) = \frac{1}{N_{\mathrm{gt}}} \sum_{i=1}^{N_{\mathrm{gt}}} \mathbf{1}\!\left[\max_{k} J(\hat{\mathbf{w}}^{(k)}, \mathbf{w}_i) \geq \tau\right],
\end{equation}
where $J$ is the way-set Jaccard similarity.
We sweep $\tau \in [0.05, 0.95]$ with step 0.01, and compute the area under the curve via trapezoidal integration.
The final CovAUC is the mean across all OD groups with sufficient ground-truth multiplicity ($N_{\mathrm{gt}} \geq 3$).

\section{Additional Ablation Details}
\label{app:ablation}

\paragraph{$D = 128$ failure analysis.}
At $D = 128$, the $\beta$-VAE KL divergence increases beyond the flow model's capacity to learn sharp conditional distributions.
The flow validation loss does not converge to a comparable level as $D = 64$ (0.435), and generated $\boldsymbol{\mu}$ samples tend to cluster near the unconditional mean rather than spread across corridor-specific regions.
This manifests as repeated generation of a single ``average'' corridor per OD pair, explaining the 24.8pp arrival rate drop.

\paragraph{Anti-loop sensitivity.}
We tested penalty values $\alpha \in \{1.0, 1.5, 2.0, 3.0\}$ and lookback windows $k \in \{2, 4, 6, 8\}$.
The combination $(\alpha = 2.0, k = 4)$ yields the best arrival rate.
Larger penalties ($\alpha = 3.0$) can suppress valid revisitations of ways that appear in both the outbound and return portions of a route; smaller lookback windows ($k = 2$) miss longer-period oscillations.

\section{Coverage-$\tau$ Detailed Results}
\label{app:cov-detail}

\begin{table}[h]
  \caption{Coverage at selected $\tau$ thresholds.}
  \label{tab:cov-detail}
  \centering
  \small
  \begin{tabular}{lcccc}
    \toprule
    Method & Cov@$\tau$=0.3 & Cov@$\tau$=0.5 & AUC \\
    \midrule
    CascadeTraj    & 0.385 & 0.137 & 0.216 \\
    Transformer AR & 0.081 & 0.020 & 0.052 \\
    RNN AR         & 0.077 & 0.023 & 0.045 \\
    Shortest Path  & 0.000 & 0.000 & 0.023 \\
    \bottomrule
  \end{tabular}
\end{table}
